\chapter{Query Rewriting and Evaluation Methodology}

Student questions often use informal or conversational language, while official university documents use administrative terminology. Even when the intent is clear to a human reader, the query may lack the lexical signals needed by BM25. The rewrite layer is therefore used as a lightweight semantic adapter before deterministic retrieval.

\section{Two-Layer Rewrite Task}
For each student question, the LLM produces two fields:

\begin{itemize}
    \item \texttt{cleaned\_original\_query}: a compact version of the original wording that preserves the information need and its key terms;
    \item \texttt{augmentation\_query}: a small number of additional lexical terms or fixed phrases likely to occur in official titles, headings, FAQs, or regulations.
\end{itemize}

The cleaned layer remains the retrieval anchor. The augmentation layer supports it without replacing the original intent. The final rewritten condition searches both layers, with the cleaned layer weighted more strongly.

The rewrite prompt forbids answering the question, inventing facts, adding unspecified course names, or expanding toward only broadly related administrative topics. This restriction is necessary because a longer query is not automatically better: unrelated but domain-plausible words can damage BM25 ranking.

\section{Experimental Conditions}
The experiment crosses two query conditions with two corpus conditions:

\begin{table}[h]
\centering
\begin{tabular}{lll}
\toprule
\textbf{Dimension} & \textbf{Condition} & \textbf{Meaning} \\
\midrule
Query & Original & Student question submitted directly to BM25 \\
Query & Rewritten & Frozen two-layer LLM rewrite submitted to BM25 \\
Corpus & Polished & Curated source set with limited unrelated material \\
Corpus & Noisy & Broader source set containing the polished documents and additional sources \\
\bottomrule
\end{tabular}
\caption{Experimental factors.}
\label{tab:experimental-factors}
\end{table}

The benchmark contains 26 queries. For each corpus, the original-query baseline is run once. Three rewrite sets are generated independently, frozen, and reused unchanged on both corpora. This produces a controlled comparison in which rewrite variability and corpus noise can be analysed separately.

All conditions use the same BM25 profile, field weights, document-collapse policy, and top-$k$ depth. Only the query representation and exact source-set filter change.

\section{Manual Relevance Judgement}
Each of the top three retrieved documents is manually assigned one of the following labels:

\begin{table}[h]
\centering
\begin{tabular}{clp{0.58\linewidth}}
\toprule
\textbf{Label} & \textbf{Name} & \textbf{Meaning} \\
\midrule
2 & Relevant & Directly provides the information needed to answer the question. \\
1 & Partially relevant & Useful but incomplete, indirect, ambiguous, or insufficient by itself. \\
0 & Irrelevant & Does not contribute useful evidence for the question. \\
\bottomrule
\end{tabular}
\caption{Manual relevance scale.}
\label{tab:relevance-scale}
\end{table}

An \emph{unclear} state is available during annotation, but unresolved hits are not accepted in the final evaluation package. Judgements follow the query--result pair: the same retrieved segment for the same case must receive the same label across runs.

Two evaluation interpretations are reported:

\begin{itemize}
    \item \textbf{Strict}: only relevant results count as useful;
    \item \textbf{Lenient}: relevant and partially relevant results both count as useful.
\end{itemize}

\section{Primary Retrieval Metrics}
Because only the top three distinct documents are retrieved and judged, the primary cutoffs are 1 and 3. Hit@5 is not reported as a separate result.

\subsection{Hit@$k$}
For query $q$, Hit@$k$ is one when at least one useful result appears within ranks $1$ to $k$, and zero otherwise:

\begin{equation}
\operatorname{Hit@k}(q)=
\begin{cases}
1, & \text{if at least one useful result occurs at rank } \leq k,\\
0, & \text{otherwise.}
\end{cases}
\end{equation}

The reported values are the mean over all evaluated queries for $k=1$ and $k=3$. Strict and lenient versions differ only in which relevance labels count as useful.

Hit@3 is the main answerability measure: it estimates the proportion of questions for which the retrieved context contains at least one usable official source.

\subsection{Mean Reciprocal Rank}
Mean Reciprocal Rank rewards systems that place the first useful result higher:

\begin{equation}
\operatorname{MRR}=\frac{1}{|Q|}\sum_{q \in Q}\frac{1}{r_q},
\end{equation}

where $r_q$ is the rank of the first useful result and the reciprocal rank is zero when no useful result appears in the top three. Strict MRR uses the first relevant result; lenient MRR uses the first relevant or partially relevant result.

\section{Additional Context Metrics}
Hit@3 measures whether the query is answerable from the retrieved window, but it does not distinguish one useful source from several complementary sources. Extra-context metrics therefore evaluate only the evidence beyond the first useful result.

For each answerable case, the \textbf{primary useful hit} is selected as follows:

\begin{enumerate}
    \item the first relevant result, when one exists;
    \item otherwise, the first partially relevant result.
\end{enumerate}

The primary hit is removed from extra-context scoring. Each remaining relevant hit contributes $1$, each remaining partial hit contributes $0.5$, and all other hits contribute $0$.

Let $A$ be the set of answerable cases, $E \subseteq A$ the cases with positive extra-context gain, and $g_q$ the extra gain for case $q$. The reported metrics are:

\begin{align}
\operatorname{ExtraCoverage@3} &= \frac{|E|}{|A|},\\
\operatorname{ExtraGainPerAnswerable@3} &= \frac{\sum_{q\in A}g_q}{|A|},\\
\operatorname{ExtraGainWhenPresent@3} &= \frac{\sum_{q\in E}g_q}{|E|}.
\end{align}

The first metric shows how often answerable queries receive additional useful context. The second keeps answerable cases with no extra evidence in the denominator. The third describes the strength of the added context only when it is present.

\section{Paired Comparison and Reproducibility}
Original and rewritten results are compared case by case. Aggregate deltas are considered valid only when both runs use the same query set, corpus source set, retrieval profile, and top-$k$ depth.

The evaluation package stores the raw retrieval runs, frozen rewrite snapshots, manual judgements, and paired comparison exports. The paired JSON export contains strict and lenient Hit@1, Hit@3, MRR, extra-context metrics, aggregate deltas, and per-case rank changes.

This design keeps the evaluated retrieval process inspectable: the LLM generates the rewrite once, while BM25 ranking, corpus filtering, manual judgements, and metric computation remain deterministic for the frozen inputs.
\par
